\chapter{Il Newsvendor e le sue varianti}
\label{cap:newsvendor}

\noindent\textbf{Classe:} convesso 1D / LP stocastico a scenari \hfill
\textbf{Script:} \texttt{python/lab12\_newsvendor.py}

\section{Motivazione gestionale}

Il Newsvendor (``venditore di giornali'') descrive una decisione comunissima:
scegliere una quantità \emph{prima} di osservare la domanda. Ordinare troppo genera
invenduto; troppo poco, vendite perse. Si applica a moda e stagionali, freschi e
deperibili, farmaci, capacità alberghiera, componenti di ricambio. È la porta
d'ingresso dell'ottimizzazione stocastica, e il suo punto di forza didattico è un
risultato netto: \textbf{la quantità ottima non è la domanda media}.



Il modello ha più di un secolo e resta il fondamento della gestione delle scorte
per prodotti a ciclo singolo. La sua importanza didattica va oltre: è il primo
incontro con l'idea che \emph{si decide prima e si scopre poi}, e con gli strumenti
per farlo bene (scenari, valore della soluzione stocastica, misure di rischio).

\textbf{In questo capitolo impareremo a:} (1) ricavare e applicare la regola del
quantile; (2) trasformare un problema stocastico in un LP a scenari con il trucco
delle parti positive; (3) quantificare il valore di modellare l'incertezza (VSS);
(4) incorporare livelli di servizio e avversione al rischio.

\section{Il modello base}

\paragraph{Il problema a parole.} \emph{Decidiamo} una sola quantità $q$, prima di
conoscere la domanda. \emph{L'obiettivo} è minimizzare il costo atteso complessivo
dell'errore: quello di aver ordinato troppo (invenduto) più quello di aver ordinato
troppo poco (vendite perse). \emph{Vincoli} non ce ne sono nel modello base ($q \ge 0$);
compaiono nelle varianti: budget condiviso, livelli di servizio, limiti di capacità.

\subsection*{Dati (input del modello)}
\begin{center}\small
\begin{tabular}{llp{7.2cm}}
\toprule
Simbolo & Tipo & Significato \\
\midrule
$D$ & variabile aleatoria $\ge 0$ & domanda, con funzione di ripartizione $F$ \\
$c_u$ & $\in \R_{> 0}$ & costo unitario di \emph{underage} (domanda non servita):
        margine perso più eventuale penalità \\
$c_o$ & $\in \R_{> 0}$ & costo unitario di \emph{overage} (unità invenduta):
        costo meno valore di recupero \\
\bottomrule
\end{tabular}
\end{center}

\subsection*{Variabile decisionale}
\begin{center}\small
\begin{tabular}{llp{7.2cm}}
\toprule
$q$ & $\ge 0$, continua & quantità ordinata \emph{prima} di osservare $D$ \\
\bottomrule
\end{tabular}
\end{center}

\begin{modello}[Newsvendor in forma di costo]
\[
\min_{q \ge 0}\;\; C(q) = c_o\, \E\bigl[(q - D)^+\bigr] + c_u\, \E\bigl[(D - q)^+\bigr]
\]
Con prezzo $p$, costo $c$ e recupero $v \le c$:\; $c_u = p - c$,\; $c_o = c - v$.
$C(q)$ è convessa; se $F$ è continua l'ottimo soddisfa la \textbf{regola del quantile}:
\[
F(q^*) = \alpha^* = \frac{c_u}{c_u + c_o}
\qquad\Longrightarrow\qquad
q^* = F^{-1}(\alpha^*).
\]
\end{modello}

\begin{spiegazione}[Da dove viene la regola del quantile]
Si ragiona sull'unità marginale: ordinare la $q$-esima unità in più rende $c_u$ se la
domanda la assorbe (probabilità $1 - F(q)$) e costa $c_o$ se resta invenduta
(probabilità $F(q)$). Conviene finché
$c_u\,(1 - F(q)) \ge c_o\,F(q)$, cioè finché $F(q) \le c_u/(c_u + c_o)$.
Se la carenza costa molto più dell'eccedenza, $\alpha^*$ è alto e si ordina un
quantile alto; se l'invenduto è costoso (deperibili), si ordina poco.
\end{spiegazione}

\section{Esempio svolto a mano (domanda discreta)}

\begin{esempio}[Tre scenari equiprobabili]
$D \in \{80, 100, 120\}$ con probabilità $1/3$; $c_u = 9$, $c_o = 4$
(quindi $\alpha^* = 9/13 = 0{,}6923$).

\textbf{Passo 1 — costo atteso per i tre candidati naturali.}
\begin{align*}
C(80) &= \tfrac13\,\bigl[0 + 9 \cdot 20 + 9 \cdot 40\bigr] = 180 \\
C(100) &= \tfrac13\,\bigl[4 \cdot 20 + 0 + 9 \cdot 20\bigr] = \tfrac{260}{3} = 86{,}7 \\
C(120) &= \tfrac13\,\bigl[4 \cdot 40 + 4 \cdot 20 + 0\bigr] = 80
\end{align*}

\textbf{Passo 2 — regola del quantile discreta.} Serve il più piccolo $q$ con
$F(q) \ge 0{,}6923$: $F(80) = 1/3$, $F(100) = 2/3 = 0{,}667 < 0{,}6923$,
$F(120) = 1$. Quindi $q^* = 120$: coerente con il passo 1.

\textbf{Passo 3 — lettura.} Ordinare il \emph{massimo} storico nonostante la media
sia 100: con $c_u = 9 \gg c_o = 4$, restare corti costa più del doppio che restare
lunghi. Chi ordina ``la media'' paga $86{,}7$ contro $80$: l'incertezza si gestisce,
non si ignora.
\end{esempio}

\section{Caso di studio: la regola del quantile in continuo}

Panetteria: $p = 15$, $c = 6$, $v = 2$~\euro{} $\Rightarrow c_u = 9$, $c_o = 4$;
domanda normale con $\mu = 100$, $\sigma = 20$.

\begin{lstlisting}[style=output]
alpha* = 9/13 = 0,6923
q* = F^-1(0,6923) = 110,05 unita'   (media = 100)
costo atteso in q*: 91,43 EUR    in q = media: 103,72 EUR
\end{lstlisting}

\begin{center}
\begin{tikzpicture}
\begin{axis}[labmezza, title={Il minimo è al 69° percentile},
  xlabel={quantità ordinata $q$}, ylabel={costo atteso (\euro)}]
\addplot [teal, very thick] table [col sep=comma, x=q, y=costo] {\figdat{cap12_costo}};
\addplot [grigiomedio, dashed] coordinates {(100, 80) (100, 450)};
\addplot [rossomattone, dash dot] coordinates {(110.05, 80) (110.05, 450)};
\legend{$C(q)$, media $= 100$, $q^* = 110$}
\end{axis}
\end{tikzpicture}%
\begin{tikzpicture}
\begin{axis}[labmezza, title={Stabilità della soluzione a scenari},
  xlabel={numero di scenari $S$}, ylabel={$q$ ottima}, xmode=log]
\addplot [teal, thick, mark=*, mark size=1.6pt,
  error bars/.cd, y dir=both, y explicit]
  table [col sep=comma, x=S, y=q_media, y error=q_std] {\figdat{cap12_stabilita}};
\addplot [rossomattone, dash dot, domain=10:3000] {110.05};
\legend{media $\pm$ dev.std su 30 repliche, $q^*$ teorico}
\end{axis}
\end{tikzpicture}
\captionof{figure}{A sinistra: il costo atteso $C(q)$ è convesso e il minimo cade in
$q^* = 110$ (69° percentile della domanda), non sulla media 100. A destra: la $q$
ottima stimata da $|S|$ scenari campionati, media e deviazione standard su 30
repliche --- con pochi scenari la decisione è instabile.}
\end{center}

\section{La formulazione lineare a scenari}

Quando $F$ non è nota in forma chiusa --- il caso normale in azienda --- si usano gli
scenari: dati storici, campioni da una distribuzione stimata, o scenari pesati.

\begin{modello}[Newsvendor come LP]
\begin{align}
\min\;& c_o \sum_{s \in S} \pi_s\, o_s + c_u \sum_{s \in S} \pi_s\, u_s \notag\\
\text{s.t.}\;\; & o_s \ge q - d_s,\qquad u_s \ge d_s - q && \forall s \in S \notag\\
& q,\, o_s,\, u_s \ge 0 \notag
\end{align}
\end{modello}

\begin{spiegazione}[Obiettivo e vincoli, riga per riga]
\textbf{Funzione obiettivo.} La media pesata, sugli scenari, dei costi di eccedenza
($c_o\, o_s$) e carenza ($c_u\, u_s$): è la versione ``a scenari'' del costo atteso
del modello base.

\textbf{Vincoli ($o_s \ge q - d_s$ e $u_s \ge d_s - q$).} Definiscono eccedenza e
carenza di ogni scenario a partire dall'unica vera decisione $q$; con i vincoli di
non negatività realizzano le parti positive, come spiegato sotto.
\end{spiegazione}

\begin{spiegazione}[Perché è lineare (il trucco delle parti positive)]
$(q - d_s)^+$ non è lineare, ma all'\emph{ottimo} la variabile $o_s$, che compare
nell'obiettivo con coefficiente positivo, si schiaccia sul massimo tra $q - d_s$ e
$0$: il vincolo $o_s \ge q - d_s$ insieme a $o_s \ge 0$ fa esattamente il lavoro di
$\max\{q - d_s, 0\}$ senza variabili binarie. Stesso meccanismo per $u_s$. Questo
trucco --- linearizzare le parti positive che l'obiettivo spinge verso il basso ---
torna identico nel CVaR (capitolo~\ref{cap:varcvar}) e nella SVM
(capitolo~\ref{cap:svm}).
\end{spiegazione}

\begin{lstlisting}[style=python]
q = m.addVar(name="q")
o = m.addVars(S, name="o")     # eccedenza per scenario
u = m.addVars(S, name="u")     # carenza per scenario
m.addConstrs((o[s] >= q - dom[s] for s in range(S)))
m.addConstrs((u[s] >= dom[s] - q for s in range(S)))
m.setObjective(gp.quicksum((Co*o[s] + Cu*u[s]) / S
                           for s in range(S)), GRB.MINIMIZE)
\end{lstlisting}

\begin{lstlisting}[style=output]
LP con S = 600 scenari: q = 108,61
quantile empirico al 69,23%: 108,60   (coincidono: e' un teorema)
VSS: decisione q = E[D] = 100 costa 99,50; q stocastica costa 88,33
     valore della soluzione stocastica = 11,17 EUR per ciclo
\end{lstlisting}

\begin{insight}
Il \emph{valore della soluzione stocastica} (VSS) risponde alla domanda del capo:
``perché complicarsi la vita con gli scenari?''. Qui vale 11,17~\euro{} a ciclo,
l'11\% del costo: è il guadagno di \emph{modellare} l'incertezza invece di sostituirla
con la media. La figura sulla stabilità dà l'altra metà del messaggio: sotto i 100
scenari la $q$ ottima balla di $\pm 3$ unità tra una stima e l'altra --- gli scenari
sono anch'essi un campione, e la decisione eredita la loro varianza.
\end{insight}

\section{Livelli di servizio}

\begin{lstlisting}[style=output]
cycle service level 90%: q = 125,63  (costo +23,41 EUR sull'ottimo economico)
cycle service level 95%: q = 132,90  (costo +45,59)
cycle service level 99%: q = 146,53  (costo +95,56)
fill rate in q*: 96,06%   probabilita' di non-stockout in q*: 69,23%
\end{lstlisting}

\begin{attenzione}[Non confondere i livelli di servizio]
In $q^* = 110$ la probabilità di evitare lo stock-out è solo il 69\%, ma il
\emph{fill rate} (quota di domanda servita in media) è il 96\%: misure diverse che i
contratti spesso confondono. Promettere ``95\% di probabilità di copertura totale''
costa $+45{,}59$~\euro{} a ciclo rispetto all'ottimo economico; promettere ``95\% di
fill rate'' è quasi gratis. Leggere bene lo SLA prima di firmarlo.
\end{attenzione}

\section{Avversione al rischio: media-CVaR}

Il costo atteso non distingue tra tante piccole perdite e rari disastri. Con la
formulazione lineare del CVaR (anticipata qui, derivata nel
capitolo~\ref{cap:varcvar}) minimizziamo
$(1-\lambda)\,\E[\text{costo}] + \lambda\, \cvar_{0,90}(\text{costo})$:

\begin{lstlisting}[style=output]
lambda = 0,00: q = 108,61  costo medio 88,33  CVaR90 245,93
lambda = 0,25: q = 112,45  costo medio 90,37  CVaR90 229,29
lambda = 0,50: q = 114,48  costo medio 92,52  CVaR90 225,49
lambda = 1,00: q = 116,17  costo medio 94,87  CVaR90 224,56
\end{lstlisting}

\begin{center}
\begin{tikzpicture}
\begin{axis}[lab, title={Frontiera costo medio -- rischio di coda},
  xlabel={CVaR$_{0{,}90}$ del costo (\euro)}, ylabel={costo medio (\euro)}]
\addplot+ [mark=*, mark size=1.8pt] table [col sep=comma, x=cvar, y=costo_medio]
  {\figdat{cap12_frontiera_cvar}};
\node [font=\scriptsize, anchor=west] at (axis cs:245.93,88.33) {$\lambda=0$};
\node [font=\scriptsize, anchor=south west] at (axis cs:224.56,94.87) {$\lambda=1$};
\end{axis}
\end{tikzpicture}
\captionof{figure}{Frontiera tra costo medio e CVaR$_{0{,}90}$ del costo, ottenuta
variando il peso $\lambda$ dell'avversione al rischio: da $\lambda = 0$ (ottimo in
media, coda pesante) a $\lambda = 1$ (protezione massima). Il primo tratto è quasi
verticale: molta protezione per pochissimo costo medio.}
\end{center}

Il decisore avverso al rischio ordina di più (da 108,6 a 116,2 unità): paga
$6{,}5$~\euro{} in più in media per tagliare di $21{,}4$~\euro{} il costo medio dei
peggiori 10\% degli scenari. La gran parte della protezione si ottiene già con
$\lambda = 0{,}25$: la frontiera è ripida all'inizio e piatta poi.

\section{Multiprodotto con budget condiviso}

Tre dolci con domande \emph{correlate} ($\rho = 0{,}7$: le feste vanno bene o male per
tutti insieme) e un budget di produzione di 1200~\euro:

\begin{lstlisting}[style=output]
  prodotto  | q senza budget | q con budget
 panettone  |          111,1 |         99,3
 pandoro    |           91,5 |         75,4
 torrone    |           65,9 |         56,8
Spesa: 1200/1200   prezzo ombra del budget: -0,553
\end{lstlisting}

\begin{spiegazione}
Il budget vincolante comprime le tre quantità sotto i rispettivi quantili ottimi, ma
non proporzionalmente: il taglio dipende dal rapporto $c_u/c_o$ e dal costo unitario
di ciascun prodotto. Il duale dice che un euro di budget in più ridurrebbe il costo
atteso di $0{,}55$~\euro: un rendimento del 55\% che giustificherebbe quasi qualunque
finanziamento --- informazione che senza il modello nessuno avrebbe quantificato.
La correlazione alta, inoltre, toglie il beneficio di diversificazione: i tre
prodotti falliscono insieme, e trattarli come indipendenti sottostimerebbe il rischio.
\end{spiegazione}

\section{Esercizi}

\begin{esercizio}[verifica]
Ricalcolare a mano $\alpha^*$ e $q^*$ se il Comune impone una penalità
$b = 3$~\euro{} per cliente non servito ($c_u = p - c + b$). Verificare con lo script.
\end{esercizio}

\begin{esercizio}[dati storici]
Usare come scenari le 24 osservazioni di domanda degli ultimi due anni (generarle con
lo script) invece della distribuzione: quanto differisce $q$? Quale approccio
consigliereste a un'azienda che ha solo 24 osservazioni?
\end{esercizio}

\begin{esercizio}[deperibili]
Il recupero diventa negativo ($v = -1$~\euro: smaltire costa): come cambiano
$\alpha^*$ e $q^*$? Tracciare $q^*$ in funzione di $v \in [-2, 4]$.
\end{esercizio}

\begin{esercizio}[EVPI]
Calcolare l'EVPI (\emph{Expected Value of Perfect Information}, valore atteso
dell'informazione perfetta): costo medio con $q$ scelta
\emph{dopo} aver visto la domanda (cioè $q = d_s$, costo 0) contro costo della
soluzione stocastica. Quanto vale al massimo una previsione perfetta della domanda?
\end{esercizio}

\begin{esercizio}[fill rate vincolato]
Aggiungere all'LP il vincolo di fill rate
$\sum_s \pi_s u_s \le (1 - 0{,}98)\,\E[D]$ e confrontare quantità e costo con il
cycle service level al 98\%.
\end{esercizio}

\begin{esercizio}[multiprodotto]
Rieseguire il multiprodotto con domande indipendenti ($\rho = 0$) a parità di tutto
il resto: il costo atteso ottimo migliora o peggiora? Perché?
\end{esercizio}
